% Options for packages loaded elsewhere
\PassOptionsToPackage{unicode}{hyperref}
\PassOptionsToPackage{hyphens}{url}
%
\documentclass[
  ignorenonframetext,
]{beamer}
\usepackage{pgfpages}
\setbeamertemplate{caption}[numbered]
\setbeamertemplate{caption label separator}{: }
\setbeamercolor{caption name}{fg=normal text.fg}
\beamertemplatenavigationsymbolsempty
% Prevent slide breaks in the middle of a paragraph
\widowpenalties 1 10000
\raggedbottom
\setbeamertemplate{part page}{
  \centering
  \begin{beamercolorbox}[sep=16pt,center]{part title}
    \usebeamerfont{part title}\insertpart\par
  \end{beamercolorbox}
}
\setbeamertemplate{section page}{
  \centering
  \begin{beamercolorbox}[sep=12pt,center]{part title}
    \usebeamerfont{section title}\insertsection\par
  \end{beamercolorbox}
}
\setbeamertemplate{subsection page}{
  \centering
  \begin{beamercolorbox}[sep=8pt,center]{part title}
    \usebeamerfont{subsection title}\insertsubsection\par
  \end{beamercolorbox}
}
\AtBeginPart{
  \frame{\partpage}
}
\AtBeginSection{
  \ifbibliography
  \else
    \frame{\sectionpage}
  \fi
}
\AtBeginSubsection{
  \frame{\subsectionpage}
}
\usepackage{amsmath,amssymb}
\usepackage{lmodern}
\usepackage{iftex}
\ifPDFTeX
  \usepackage[T1]{fontenc}
  \usepackage[utf8]{inputenc}
  \usepackage{textcomp} % provide euro and other symbols
\else % if luatex or xetex
  \usepackage{unicode-math}
  \defaultfontfeatures{Scale=MatchLowercase}
  \defaultfontfeatures[\rmfamily]{Ligatures=TeX,Scale=1}
\fi
\usetheme[]{Montpellier}
% Use upquote if available, for straight quotes in verbatim environments
\IfFileExists{upquote.sty}{\usepackage{upquote}}{}
\IfFileExists{microtype.sty}{% use microtype if available
  \usepackage[]{microtype}
  \UseMicrotypeSet[protrusion]{basicmath} % disable protrusion for tt fonts
}{}
\makeatletter
\@ifundefined{KOMAClassName}{% if non-KOMA class
  \IfFileExists{parskip.sty}{%
    \usepackage{parskip}
  }{% else
    \setlength{\parindent}{0pt}
    \setlength{\parskip}{6pt plus 2pt minus 1pt}}
}{% if KOMA class
  \KOMAoptions{parskip=half}}
\makeatother
\usepackage{xcolor}
\IfFileExists{xurl.sty}{\usepackage{xurl}}{} % add URL line breaks if available
\IfFileExists{bookmark.sty}{\usepackage{bookmark}}{\usepackage{hyperref}}
\hypersetup{
  hidelinks,
  pdfcreator={LaTeX via pandoc}}
\urlstyle{same} % disable monospaced font for URLs
\newif\ifbibliography
\usepackage{color}
\usepackage{fancyvrb}
\newcommand{\VerbBar}{|}
\newcommand{\VERB}{\Verb[commandchars=\\\{\}]}
\DefineVerbatimEnvironment{Highlighting}{Verbatim}{commandchars=\\\{\}}
% Add ',fontsize=\small' for more characters per line
\usepackage{framed}
\definecolor{shadecolor}{RGB}{248,248,248}
\newenvironment{Shaded}{\begin{snugshade}}{\end{snugshade}}
\newcommand{\AlertTok}[1]{\textcolor[rgb]{0.94,0.16,0.16}{#1}}
\newcommand{\AnnotationTok}[1]{\textcolor[rgb]{0.56,0.35,0.01}{\textbf{\textit{#1}}}}
\newcommand{\AttributeTok}[1]{\textcolor[rgb]{0.77,0.63,0.00}{#1}}
\newcommand{\BaseNTok}[1]{\textcolor[rgb]{0.00,0.00,0.81}{#1}}
\newcommand{\BuiltInTok}[1]{#1}
\newcommand{\CharTok}[1]{\textcolor[rgb]{0.31,0.60,0.02}{#1}}
\newcommand{\CommentTok}[1]{\textcolor[rgb]{0.56,0.35,0.01}{\textit{#1}}}
\newcommand{\CommentVarTok}[1]{\textcolor[rgb]{0.56,0.35,0.01}{\textbf{\textit{#1}}}}
\newcommand{\ConstantTok}[1]{\textcolor[rgb]{0.00,0.00,0.00}{#1}}
\newcommand{\ControlFlowTok}[1]{\textcolor[rgb]{0.13,0.29,0.53}{\textbf{#1}}}
\newcommand{\DataTypeTok}[1]{\textcolor[rgb]{0.13,0.29,0.53}{#1}}
\newcommand{\DecValTok}[1]{\textcolor[rgb]{0.00,0.00,0.81}{#1}}
\newcommand{\DocumentationTok}[1]{\textcolor[rgb]{0.56,0.35,0.01}{\textbf{\textit{#1}}}}
\newcommand{\ErrorTok}[1]{\textcolor[rgb]{0.64,0.00,0.00}{\textbf{#1}}}
\newcommand{\ExtensionTok}[1]{#1}
\newcommand{\FloatTok}[1]{\textcolor[rgb]{0.00,0.00,0.81}{#1}}
\newcommand{\FunctionTok}[1]{\textcolor[rgb]{0.00,0.00,0.00}{#1}}
\newcommand{\ImportTok}[1]{#1}
\newcommand{\InformationTok}[1]{\textcolor[rgb]{0.56,0.35,0.01}{\textbf{\textit{#1}}}}
\newcommand{\KeywordTok}[1]{\textcolor[rgb]{0.13,0.29,0.53}{\textbf{#1}}}
\newcommand{\NormalTok}[1]{#1}
\newcommand{\OperatorTok}[1]{\textcolor[rgb]{0.81,0.36,0.00}{\textbf{#1}}}
\newcommand{\OtherTok}[1]{\textcolor[rgb]{0.56,0.35,0.01}{#1}}
\newcommand{\PreprocessorTok}[1]{\textcolor[rgb]{0.56,0.35,0.01}{\textit{#1}}}
\newcommand{\RegionMarkerTok}[1]{#1}
\newcommand{\SpecialCharTok}[1]{\textcolor[rgb]{0.00,0.00,0.00}{#1}}
\newcommand{\SpecialStringTok}[1]{\textcolor[rgb]{0.31,0.60,0.02}{#1}}
\newcommand{\StringTok}[1]{\textcolor[rgb]{0.31,0.60,0.02}{#1}}
\newcommand{\VariableTok}[1]{\textcolor[rgb]{0.00,0.00,0.00}{#1}}
\newcommand{\VerbatimStringTok}[1]{\textcolor[rgb]{0.31,0.60,0.02}{#1}}
\newcommand{\WarningTok}[1]{\textcolor[rgb]{0.56,0.35,0.01}{\textbf{\textit{#1}}}}
\setlength{\emergencystretch}{3em} % prevent overfull lines
\providecommand{\tightlist}{%
  \setlength{\itemsep}{0pt}\setlength{\parskip}{0pt}}
\setcounter{secnumdepth}{-\maxdimen} % remove section numbering
%\documentclass[unicode,10pt]{beamer}% 'unicode'が必要
\usepackage{luatexja}% 日本語を使う場合は必須
%\usepackage[japanese]{babel}	
%\usefonttheme[onlymath]{serif}
%\usepackage[T1]{fontenc} %これを使うと、ギリシャ文字が出ない
%\usepackage{textcomp}
%\usepackage[scale=1.0]{tgheros} % Sans serif font
%\usepackage[scaled]{beramono} % Monospace font
%\usepackage{luatexja-otf}
%\renewcommand{\kanjifamilydefault}{\gtdefault}
%\usepackage[haranoaji]{luatexja-preset}

% スライド番号の表示 %%%%%%%%%%%%%%%%%%%
\makeatletter
\setbeamertemplate{footline}{%
  \leavevmode%
  \hbox{%
  \begin{beamercolorbox}[wd=.5\paperwidth,ht=2.25ex,dp=1ex,right]{date in head/foot}%
    \insertframenumber{} \hspace*{2ex} % new version without total frames
  \end{beamercolorbox}}%
  \vskip0pt%
}
\makeatother
% スライド番号の表示 %%%%%%%%%%%%%%%%%%%
\ifLuaTeX
  \usepackage{selnolig}  % disable illegal ligatures
\fi

\title{第2章: 因果関係\\
(2.6. 1変数の記述統計量)}
\subtitle{今井耕介 著\\
『社会科学のためのデータ分析入門 (QSS)』}
\author{}
\date{\vspace{-2.5em}2026-03-09}

\begin{document}
\frame{\titlepage}

\hypertarget{ux5909ux6570ux306eux8a18ux8ff0ux7d71ux8a08ux91cf}{%
\section{2.6
1変数の記述統計量}\label{ux5909ux6570ux306eux8a18ux8ff0ux7d71ux8a08ux91cf}}

\begin{frame}{記述統計の目的}
\protect\hypertarget{ux8a18ux8ff0ux7d71ux8a08ux306eux76eeux7684}{}
\begin{itemize}
\tightlist
\item
  大規模なデータの分布（散らばり具合や中心的な傾向）を、少数の数値で要約する。
\item
  主要な指標：

  \begin{enumerate}
  \tightlist
  \item
    \textbf{中心傾向の尺度}: 平均値 (Mean)、中央値 (Median)
  \item
    \textbf{散らばりの尺度}: 分位数 (Quantiles)、四分位範囲
    (IQR)、標準偏差 (Standard Deviation)
  \end{enumerate}
\end{itemize}
\end{frame}

\hypertarget{ux5206ux4f4dux6570-quantiles}{%
\section{2.6.1 分位数 (Quantiles)}\label{ux5206ux4f4dux6570-quantiles}}

\begin{frame}[fragile]{minwage データの読み込み (1)}
\protect\hypertarget{minwage-ux30c7ux30fcux30bfux306eux8aadux307fux8fbcux307f-1}{}
\begin{itemize}
\tightlist
\item
  最低賃金データの読み込みと準備を行います。
\end{itemize}

\scriptsize

\begin{Shaded}
\begin{Highlighting}[]
\CommentTok{\# 1. ローカルに保存したminwageデータの読み込み（推奨）}
\NormalTok{minwage }\OtherTok{\textless{}{-}} \FunctionTok{read.csv}\NormalTok{(}\StringTok{"minwage.csv"}\NormalTok{)}

\CommentTok{\# （参考）URLから直接読み込むことも可能}
\CommentTok{\# minwage \textless{}{-} read.csv("https://ayumu{-}tanaka.github.io/QSS/QSS\_Data/minwage.csv")}

\CommentTok{\# 2. データの準備（州別の抽出）}
\NormalTok{minwageNJ }\OtherTok{\textless{}{-}} \FunctionTok{subset}\NormalTok{(minwage, }\AttributeTok{subset =}\NormalTok{ (location }\SpecialCharTok{!=} \StringTok{"PA"}\NormalTok{))}
\NormalTok{minwagePA }\OtherTok{\textless{}{-}} \FunctionTok{subset}\NormalTok{(minwage, }\AttributeTok{subset =}\NormalTok{ (location }\SpecialCharTok{==} \StringTok{"PA"}\NormalTok{))}
\end{Highlighting}
\end{Shaded}

\normalsize
\end{frame}

\begin{frame}[fragile]{中央値 (Median) (2)}
\protect\hypertarget{ux4e2dux592eux5024-median-2}{}
\begin{itemize}
\tightlist
\item
  データを大きさの順に並べたときに、ちょうど真ん中にくる値。
\item
  外れ値（極端に大きい・小さい値）の影響を受けにくい（ロバスト）。
\end{itemize}

\scriptsize

\begin{Shaded}
\begin{Highlighting}[]
\CommentTok{\# NJ州の引き上げ後のフルタイム雇用率の中央値を計算}
\FunctionTok{median}\NormalTok{(minwageNJ}\SpecialCharTok{$}\NormalTok{fullAfter }\SpecialCharTok{/}\NormalTok{ (minwageNJ}\SpecialCharTok{$}\NormalTok{fullAfter }\SpecialCharTok{+}\NormalTok{ minwageNJ}\SpecialCharTok{$}\NormalTok{partAfter), }
       \AttributeTok{na.rm =} \ConstantTok{TRUE}\NormalTok{)}
\end{Highlighting}
\end{Shaded}

\begin{verbatim}
## [1] 0.2916667
\end{verbatim}

\normalsize
\end{frame}

\begin{frame}[fragile]{4分位数と summary() 関数}
\protect\hypertarget{ux5206ux4f4dux6570ux3068-summary-ux95a2ux6570}{}
\begin{itemize}
\tightlist
\item
  \textbf{4分位数}: データを4等分する境界の値（25\%, 50\%, 75\%）。
\item
  \texttt{summary()}
  を使うと、最小値、最大値、平均、中央値、4分位数を一括で確認できる。
\end{itemize}

\scriptsize

\begin{Shaded}
\begin{Highlighting}[]
\CommentTok{\# NJ州の引き上げ前の賃金 (wageBefore) の要約}
\FunctionTok{summary}\NormalTok{(minwageNJ}\SpecialCharTok{$}\NormalTok{wageBefore)}
\end{Highlighting}
\end{Shaded}

\begin{verbatim}
##    Min. 1st Qu.  Median    Mean 3rd Qu.    Max. 
##    4.25    4.25    4.50    4.61    4.87    5.75
\end{verbatim}

\begin{Shaded}
\begin{Highlighting}[]
\CommentTok{\# 4分位範囲 (Interquartile Range: IQR)}
\CommentTok{\# 75\%点 と 25\%点 の差（真ん中の50\%のデータの広がり）}
\FunctionTok{IQR}\NormalTok{(minwageNJ}\SpecialCharTok{$}\NormalTok{wageBefore)}
\end{Highlighting}
\end{Shaded}

\begin{verbatim}
## [1] 0.62
\end{verbatim}

\normalsize
\end{frame}

\begin{frame}[fragile]{任意の分位数の計算 (quantile関数)}
\protect\hypertarget{ux4efbux610fux306eux5206ux4f4dux6570ux306eux8a08ux7b97-quantileux95a2ux6570}{}
\begin{itemize}
\tightlist
\item
  10\%刻み（10分位数）などで分布を詳しく見る場合に便利。
\end{itemize}

\scriptsize

\begin{Shaded}
\begin{Highlighting}[]
\CommentTok{\# 0\%から100\%まで10\%刻みの分位数を計算}
\FunctionTok{quantile}\NormalTok{(minwageNJ}\SpecialCharTok{$}\NormalTok{wageBefore, }\AttributeTok{probs =} \FunctionTok{seq}\NormalTok{(}\AttributeTok{from =} \DecValTok{0}\NormalTok{, }\AttributeTok{to =} \DecValTok{1}\NormalTok{, }\AttributeTok{by =} \FloatTok{0.1}\NormalTok{))}
\end{Highlighting}
\end{Shaded}

\begin{verbatim}
##   0%  10%  20%  30%  40%  50%  60%  70%  80%  90% 100% 
## 4.25 4.25 4.25 4.25 4.50 4.50 4.65 4.75 5.00 5.00 5.75
\end{verbatim}

\normalsize
\end{frame}

\hypertarget{ux6a19ux6e96ux504fux5dee-standard-deviation}{%
\section{2.6.2 標準偏差 (Standard
Deviation)}\label{ux6a19ux6e96ux504fux5dee-standard-deviation}}

\begin{frame}[fragile]{平均からのばらつき}
\protect\hypertarget{ux5e73ux5747ux304bux3089ux306eux3070ux3089ux3064ux304d}{}
\begin{itemize}
\tightlist
\item
  \textbf{標準偏差}:
  各データが平均からどれくらい離れているかの平均的な距離。
\item
  値が大きいほど、データが平均の周りに広く散らばっていることを示す。
\end{itemize}

\scriptsize

\begin{Shaded}
\begin{Highlighting}[]
\CommentTok{\# NJ州の引き上げ前後の雇用率の変化を計算}
\NormalTok{minwageNJ}\SpecialCharTok{$}\NormalTok{fullPropBefore }\OtherTok{\textless{}{-}}\NormalTok{ minwageNJ}\SpecialCharTok{$}\NormalTok{fullBefore }\SpecialCharTok{/} 
\NormalTok{    (minwageNJ}\SpecialCharTok{$}\NormalTok{fullBefore }\SpecialCharTok{+}\NormalTok{ minwageNJ}\SpecialCharTok{$}\NormalTok{partBefore)}
\NormalTok{minwageNJ}\SpecialCharTok{$}\NormalTok{fullPropAfter }\OtherTok{\textless{}{-}}\NormalTok{ minwageNJ}\SpecialCharTok{$}\NormalTok{fullAfter }\SpecialCharTok{/} 
\NormalTok{    (minwageNJ}\SpecialCharTok{$}\NormalTok{fullAfter }\SpecialCharTok{+}\NormalTok{ minwageNJ}\SpecialCharTok{$}\NormalTok{partAfter)}
\NormalTok{minwageNJ}\SpecialCharTok{$}\NormalTok{diff }\OtherTok{\textless{}{-}}\NormalTok{ minwageNJ}\SpecialCharTok{$}\NormalTok{fullPropAfter }\SpecialCharTok{{-}}\NormalTok{ minwageNJ}\SpecialCharTok{$}\NormalTok{fullPropBefore}

\CommentTok{\# 雇用率の変化の標準偏差を計算}
\FunctionTok{sd}\NormalTok{(minwageNJ}\SpecialCharTok{$}\NormalTok{diff, }\AttributeTok{na.rm =} \ConstantTok{TRUE}\NormalTok{)}
\end{Highlighting}
\end{Shaded}

\begin{verbatim}
## [1] 0.3010377
\end{verbatim}

\normalsize
\end{frame}

\hypertarget{ux307eux3068ux3081}{%
\section{2.6.3 まとめ}\label{ux307eux3068ux3081}}

\begin{frame}[fragile]{このセクションのまとめ}
\protect\hypertarget{ux3053ux306eux30bbux30afux30b7ux30e7ux30f3ux306eux307eux3068ux3081}{}
\begin{itemize}
\tightlist
\item
  \textbf{中央値と分位数}:
  データの順位に基づく要約。外れ値に強い（ロバスト）。
\item
  \textbf{summary()}:
  データの分布を素早く把握するための最も基本的なコマンド。
\item
  \textbf{標準偏差}: 平均の周りのばらつきを測定する。
\item
  \textbf{Rの操作}:

  \begin{itemize}
  \tightlist
  \item
    \texttt{median()}, \texttt{quantile()}, \texttt{IQR()}
    を用いた分位数の計算。
  \item
    \texttt{sd()}, \texttt{var()} を用いたばらつきの計算。
  \item
    \texttt{summary()} による一括集計。
  \item
    \texttt{na.rm\ =\ TRUE} オプションによる欠損値の除外。
  \end{itemize}
\end{itemize}
\end{frame}

\end{document}
